%! Unit 1: Algebraic Expressions and Factoring — Summative Assessment — Unit Test (blank)
%
% PARALLEL FORM of tests/practice_test/main.tex: same blueprint (A 8x1, B 6x2,
% C 10x5 with one item per lesson in lesson order, D 2x15), same difficulty,
% DIFFERENT numbers and contexts, and the vocabulary definitions are reworded and
% relettered so the practice answer letters transfer nothing. No "this is a
% practice test" remindbox — this form is distributed at test time only and is
% never merged into any packet.
%
% WORK RULE: every computation is followed by a \begin{work} block authored
% byte-identically here and in ../../test_keys/actual_test_key/main.tex.
% \workrowsep is the handwriting room and moves both files together; 8pt is the
% largest value that keeps this form and the practice form at four pages.
\documentclass[10pt]{article}
\usepackage{atda-article}
\usepackage{atda-boxes}
\newcommand{\TallMath}[1]{$\displaystyle #1\rule[-1.4em]{0pt}{3.2em}$}
\setlength{\workrowsep}{8pt}

% Part-divider strip.
% BOXGUARD: the guard lives inside the macro, so a part header can never be
% stranded alone at the foot of a page, in the blank or the key — both files use
% this identical definition.
\newcommand{\parthead}[1]{%
  \par\boxguard[9]%
  \vspace{6pt}\noindent
  \begin{headlinebox}{royal}{\color{white}\bfseries #1}\end{headlinebox}%
  \vspace{4pt}}

\begin{document}

\pageheader{Unit 1: Algebraic Expressions and Factoring}{Unit Test}
\namedateperiod

\vspace{0.06in}
\noindent\textbf{Instructions:} Show all work. No notes unless your teacher says so.
The test has four parts and is worth \textbf{100 points}.

% ======================================================================
\parthead{Part A --- Vocabulary \quad (8 items, 1 point each --- 8 points)}

\noindent\small
Write the letter of the best definition on the line beside each term. \textbf{Two definitions
will not be used.}

\vspace{5pt}
\noindent
\begin{tabularx}{\linewidth}{@{} p{0.95cm} X p{0.95cm} X @{}}
\blank{0.7cm} & \textbf{1.} Excluded value
  & \blank{0.7cm} & \textbf{5.} Rational expression \\[4pt]
\blank{0.7cm} & \textbf{2.} Difference of squares
  & \blank{0.7cm} & \textbf{6.} Perfect square trinomial \\[4pt]
\blank{0.7cm} & \textbf{3.} Zero product property
  & \blank{0.7cm} & \textbf{7.} Double root \\[4pt]
\blank{0.7cm} & \textbf{4.} Greatest common factor (GCF)
  & \blank{0.7cm} & \textbf{8.} Prime polynomial \\
\end{tabularx}

\vspace{6pt}
\begin{enumerate}[label=(\Alph*), leftmargin=2.4em, itemsep=1pt, topsep=2pt, font=\bfseries]
  \small
  \item A trinomial of the form $a^2 \pm 2ab + b^2$; it factors as $(a \pm b)^2$.
  \item A polynomial with integer coefficients that cannot be factored any further over the integers.
  \item The exponent that tells how many times a base is used as a factor.
  \item A binomial of the form $a^2-b^2$; it factors as $(a+b)(a-b)$.
  \item A repeated solution, produced by a factor that is squared.
  \item A fraction whose numerator and denominator are both polynomials.
  \item A value of the variable that makes a denominator zero.
  \item The largest monomial that divides every term of a polynomial.
  \item An expression written as a product of two or more expressions.
  \item If a product of factors is zero, then at least one of those factors must be zero.
\end{enumerate}

% ======================================================================
\parthead{Part B --- Multiple Choice \quad (6 items, 2 points each --- 12 points)}

\noindent\small Circle the best answer.

\begin{enumerate}[leftmargin=1.8em, itemsep=6pt, topsep=5pt]
\small

\item Which expression is equal to $(3n+2)^2$?
\begin{enumerate}[label=(\Alph*), leftmargin=2.2em, itemsep=1pt, topsep=2pt]
  \item $9n^2+4$
  \item $9n^2+6n+4$
  \item $9n^2+12n+4$
  \item $3n^2+12n+4$
\end{enumerate}

\item What is the value of $\sqrt{16+9}$?
\begin{enumerate}[label=(\Alph*), leftmargin=2.2em, itemsep=1pt, topsep=2pt]
  \item $7$
  \item $5$
  \item $25$
  \item $12$
\end{enumerate}

\item Which is the \textbf{complete} factorization of $8x^2+12x$?
\begin{enumerate}[label=(\Alph*), leftmargin=2.2em, itemsep=1pt, topsep=2pt]
  \item $2x(4x+6)$
  \item $4(2x^2+3x)$
  \item $4x(2x+3)$
  \item $4x(2x^2+3)$
\end{enumerate}

% BOXGUARD: without this, item 4 split across the page — stem plus options (A)
% and (B) at the foot of page 1, (C) and (D) alone at the top of page 2. A
% multiple-choice item must never break: a student cannot compare options they
% have to turn the page to see. 6 lines of item + itemsep, so [7]. Mirrored in
% ../../test_keys/actual_test_key/main.tex; the page count held at 4 in both.
\boxguard[7]
\item Which polynomial is \textbf{prime} over the integers?
\begin{enumerate}[label=(\Alph*), leftmargin=2.2em, itemsep=1pt, topsep=2pt]
  \item $x^2-36$
  \item $x^2+36$
  \item $x^2-12x+36$
  \item $x^2+12x+36$
\end{enumerate}

\item How many solutions does $x^3+16x=0$ have, and of what type?
\begin{enumerate}[label=(\Alph*), leftmargin=2.2em, itemsep=1pt, topsep=2pt]
  \item $3$ solutions: $3$ real
  \item $3$ solutions: $1$ real and $2$ imaginary
  \item $1$ solution: $1$ real
  \item $2$ solutions: $2$ real
\end{enumerate}

\item Which simplification is \textbf{correct}?
\begin{enumerate}[label=(\Alph*), leftmargin=2.2em, itemsep=1pt, topsep=2pt]
  \item $\dfrac{x+3}{x+8} = \dfrac{3}{8}$
  \item $\dfrac{x+3}{3} = x$
  \item $\dfrac{3(x+3)}{3(x+8)} = \dfrac{x+3}{x+8}$
  \item $\dfrac{x^2+3}{x^2} = 3$
\end{enumerate}

\end{enumerate}

% ======================================================================
\parthead{Part C --- Short Answer \& Computation \quad (10 items, 5 points each --- 50 points)}

\noindent\small Show all work. An answer with no supporting work earns no credit.

\begin{enumerate}[leftmargin=1.8em, itemsep=4pt, topsep=5pt]
\small

\item Simplify completely. Write your answer with positive exponents only.
\ \TallMath{\frac{\left(2a^3b^2\right)^3 \cdot 3ab}{4a^5b^4}}
\begin{work}
  \frac{\left(2a^3b^2\right)^3 \cdot 3ab}{4a^5b^4}
      &= \frac{8a^9b^6 \cdot 3ab}{4a^5b^4} \\
      &= \frac{24a^{10}b^7}{4a^5b^4} \\
      &= 6a^5b^3
\end{work}

\item (a) Simplify $\sqrt{98x^5}$, where $x>0$. \quad (b) Write $\sqrt[4]{x^3}$ using a rational
exponent. \quad (c) Evaluate $27^{2/3}$.
\begin{work}
  \text{(a)} \quad \sqrt{98x^5} &= \sqrt{49x^4}\cdot\sqrt{2x} = 7x^2\sqrt{2x} \\
  \text{(b)} \quad \sqrt[4]{x^3} &= x^{3/4} \\
  \text{(c)} \quad 27^{2/3} &= \left(\sqrt[3]{27}\right)^2 = 3^2 = 9
\end{work}

\item Factor completely by grouping: \ $2x^3+10x^2+3x+15$
\begin{work}
  2x^3+10x^2+3x+15 &= 2x^2(x+5) + 3(x+5) \\
      &= (x+5)\left(2x^2+3\right)
\end{work}

\item Factor completely: \ $6x^2+13x+5$
\begin{work}
  ac = 6\cdot 5 &= 30, \quad \text{and } 3+10 = 13 \\
  6x^2+13x+5 &= 6x^2+3x+10x+5 \\
      &= 3x(2x+1) + 5(2x+1) \\
      &= (2x+1)(3x+5)
\end{work}

\item Factor completely. \quad (a) $25x^2-81$ \qquad (b) $x^3+125$
\begin{work}
  \text{(a)} \quad 25x^2-81 &= (5x)^2 - 9^2 = (5x+9)(5x-9) \\
  \text{(b)} \quad x^3+125 &= x^3 + 5^3 = (x+5)\left(x^2-5x+25\right)
\end{work}

\item Solve by factoring: \ $3x^2 = 7x+6$
\begin{work}
  3x^2-7x-6 &= 0 \\
  (x-3)(3x+2) &= 0 \\
  x-3 = 0 \quad \text{or} \quad 3x+2 &= 0 \\
  x = 3 \quad \text{or} \quad x &= -\tfrac{2}{3}
\end{work}

\item Simplify, and state \textbf{every} excluded value.
\ \TallMath{\frac{x^2-9}{x^2+2x-15}}
\begin{work}
  \frac{x^2-9}{x^2+2x-15} &= \frac{(x+3)(x-3)}{(x+5)(x-3)} \\
      &= \frac{x+3}{x+5} \qquad (x \ne 3, \; x \ne -5)
\end{work}

\item Divide and simplify, and state \textbf{every} excluded value.
\ \TallMath{\frac{x^2-25}{x^2+8x+16} \div \frac{x-5}{x+4}}
\begin{work}
  \frac{(x+5)(x-5)}{(x+4)^2} \cdot \frac{x+4}{x-5}
      &= \frac{x+5}{x+4} \qquad (x \ne -4, \; x \ne 5)
\end{work}

\item A rectangular banner has area $A(x) = 3x^2+14x+8$ square feet and width $x+4$ feet.
(a) Find an expression for its length. (b) Find the length, width, and area when $x=6$.
\begin{work}
  3x^2+14x+8 &= (x+4)(3x+2) \\
  \text{(a)} \quad \text{length} &= 3x+2 \text{ feet} \\
  \text{(b) at } x=6: \quad \text{width} &= 10 \text{ ft}, \quad \text{length} = 20 \text{ ft} \\
  \text{area} &= 10 \cdot 20 = 200 \text{ square feet}
\end{work}

\item Factor completely. \quad (a) $4x^3-64x$ \qquad (b) $3x^4-48$
\begin{work}
  \text{(a)} \quad 4x^3-64x &= 4x\left(x^2-16\right) = 4x(x+4)(x-4) \\
  \text{(b)} \quad 3x^4-48 &= 3\left(x^4-16\right) \\
      &= 3\left(x^2+4\right)\left(x^2-4\right) \\
      &= 3\left(x^2+4\right)(x+2)(x-2)
\end{work}

\end{enumerate}

% ======================================================================
\parthead{Part D --- Extended Response \quad (2 items, 15 points each --- 30 points)}

\begin{enumerate}[leftmargin=1.8em, itemsep=8pt, topsep=5pt]
\small

\item \textbf{Verify, then judge a claim.}

\begin{enumerate}[label=(\alph*), leftmargin=1.8em, itemsep=4pt, topsep=4pt]
  \item Verify the identity $(x+3)\left(x^2-3x+9\right) = x^3+27$ by multiplying.
\begin{work}
  (x+3)\left(x^2-3x+9\right)
      &= x^3-3x^2+9x + 3x^2-9x+27 \\
      &= x^3+27
\end{work}

  \item A student writes: ``$\dfrac{x^3+27}{x+3}$ simplifies to $x^2-3x+9$, \emph{for every value
        of $x$}.'' Test both expressions at $x=1$ and at $x=-3$.
\begin{work}
  \text{at } x=1: \quad \frac{1+27}{4} &= 7, \qquad 1-3+9 = 7 \\
  \text{at } x=-3: \quad \frac{-27+27}{0} &= \frac{0}{0} \text{ --- undefined} \\
  \text{but} \quad 9+9+9 &= 27
\end{work}

  \item Is the student's algebra correct? Is the student's \emph{claim} correct? Explain in one or
        two sentences, using the word \textbf{excluded}.

\writelines{3}
\end{enumerate}

\item \textbf{A rocket is launched from a platform.} Its height in feet after $t$ seconds is
$h(t) = -16t^2+96t+112$.

\begin{enumerate}[label=(\alph*), leftmargin=1.8em, itemsep=4pt, topsep=4pt]
  \item Factor $h(t)$ completely and find both values of $t$ for which $h(t)=0$.
\begin{work}
  h(t) &= -16t^2+96t+112 \\
      &= -16\left(t^2-6t-7\right) \\
      &= -16(t-7)(t+1) \\
  t &= 7 \quad \text{or} \quad t = -1
\end{work}

  \item Which of those two values describes when the rocket lands? What does the other value mean
        here, and why is it rejected?

\writelines{2}

  \item Solve $h(t)=112$. Interpret \textbf{both} solutions --- this time neither one is rejected.
\begin{work}
  -16t^2+96t+112 &= 112 \\
  -16t^2+96t &= 0 \\
  -16t(t-6) &= 0 \\
  t = 0 \quad \text{or} \quad t &= 6
\end{work}

\writelines{2}

  \item How many solutions does the equation $h(t)=0$ have, and of what type? How does the degree
        of $h(t)$ tell you this before you factor anything?

\writelines{2}
\end{enumerate}

\end{enumerate}

\end{document}
